%% toc: P2 22 22.1,22.2,22.3
\section{Sort conventions and the one-sorted reduction in detail}
\label{sec:p02-w02-p02-038-many-sorted-reduction}

\Cref{sec:p02-many-sorted} fixes strict many-sorted signatures
(\cref{def:p02-many-sorted-signature}), strict structures and satisfaction
(\cref{def:p02-many-sorted-structure}), the guard translation
\(\operatorname{tr}_{\mathrm{MS}}\) with the sort theory
\(\mathsf{Ax}_{\mathrm{sort}}(\Lang_{\mathrm{MS}})\)
(\cref{def:p02-many-sorted-translation}) and the assignment guard
\(\mathsf{VarGuard}_{\mathrm{MS}}(\Lang_{\mathrm{MS}})\)
(\cref{def:p02-ms-variable-guard}), and states the reduction of strict
many-sorted consequence to one-sorted consequence
(\cref{thm:p02-many-sorted-translation}).  This section does four things.
It writes out the term-value and satisfaction clauses on which every proof
below depends, for a class of sorted structures wide enough to contain the
variants.  It states which sort conventions the reduction imposes, and in
particular that it imposes no coverage condition, and it gives the coverage
variant with its precondition.  It proves the satisfaction correspondence in
both directions, case by case, for the restriction of a one-sorted model and
for the tagged one-sorted structure built from a many-sorted one.  Finally it
presents two alternatives to the strict convention, one admitting empty sorts
and one admitting overlapping sorts, and shows how the reduction must change
for each.

\emph{Inputs.}  A strict many-sorted signature \(\Lang_{\mathrm{MS}}\) in the
sense of \cref{def:p02-many-sorted-signature}, its one-sorted flattening
\(\Lang_{\mathrm{MS}}^{\flat}\), and the one-sorted satisfaction relation of
\cref{sec:fol-semantics-02}, whose connectives are \(\bottom\), \(\neg\),
\(\land\), \(\lor\) and \(\to\).  Many-sorted formulas use the same
connectives.  \emph{Conclusions.}  Exact consequence-preserving reductions to
one-sorted first-order logic for the strict convention, for strict structures
with an added coverage axiom when there are finitely many sorts, for
structures with possibly empty sorts, and for structures with possibly
overlapping sorts and cross-sort equality; a proof that coverage cannot be
axiomatized when there are infinitely many sorts; and comparisons of each
variant with the strict convention.  \emph{Ambient assumptions.}  The
metatheory is ZFC.  The Axiom of Foundation is used to name fresh junk
points.  The Axiom of Choice is used only where a statement says so.
One-sorted structures have nonempty domains, as in
\cref{sec:fol-semantics-02}.  As in \cref{def:p02-many-sorted-translation},
each sorted variable is identified with its image under the fixed injection
into the one-sorted variable supply; one-sorted variables outside that image
never occur in a translated formula.  Such an injection exists exactly when
the sorted variables are countable.
\Cref{thm:p02-w02b-p02-033-variable-form} derives the strict correspondence
and the coverage variant below
(\cref{thm:p02-w02-p02-038-strict-correspondence,thm:p02-w02-p02-038-coverage-reduction})
for function and relation symbols of every cardinality, and
\cref{sec:p02-w02b-p02-033-any-cardinality} gives the reduction for
uncountably many sorts, where no such injection exists
(\cref{thm:p02-w02b-p02-033-no-renaming}).

\subsection{Sorted structures, contexts and satisfaction}
\label{sec:p02-w02-p02-038-clauses}

Write \(\mathrm{Var}_{\mathrm{MS}}\) for the set of all sorted variables of
\(\Lang_{\mathrm{MS}}\), and \(\FV(\varphi)\) for the set of sorted variables
free in a many-sorted formula \(\varphi\), defined by the usual recursion.
For each sort \(\mathsf s\), the variable written \(x^{\mathsf s}\) in
\cref{def:p02-many-sorted-signature} is a fixed member of the stock of
\(\mathsf s\); it is used below as the designated variable of that sort.

\begin{hbdefinition}{def:p02-w02-p02-038-sorted-structure}{Sorted structures, contexts and contextual assignments}
A \emph{sorted \(\Lang_{\mathrm{MS}}\)-structure} \(\mathcal M\) consists of
a set \(M_{\mathsf s}\) for each \(\mathsf s\in\mathsf{Sort}\), possibly
empty and not required to be disjoint from the others, together with
\[
\begin{gathered}
 c^{\mathcal M}\in M_{\mathsf s}\quad(c:\longrightarrow\mathsf s),\\
 f^{\mathcal M}:M_{\mathsf s_1}\times\cdots\times M_{\mathsf s_n}
   \longrightarrow M_{\mathsf s},\qquad
 P^{\mathcal M}\subseteq M_{\mathsf s_1}\times\cdots\times M_{\mathsf s_n},
\end{gathered}
\]
for the declarations
\(f:\mathsf s_1\times\cdots\times\mathsf s_n\longrightarrow\mathsf s\) with
\(n\geq1\) and \(P:\mathsf s_1\times\cdots\times\mathsf s_n\).  When the
product is empty, \(f^{\mathcal M}\) is the empty function.  The structure is
\begin{itemize}
\item \emph{strict} if every \(M_{\mathsf s}\) is nonempty and
      \(M_{\mathsf s}\cap M_{\mathsf t}=\varnothing\) for
      \(\mathsf s\ne\mathsf t\) (these are the structures of
      \cref{def:p02-many-sorted-structure});
\item an \emph{E-structure} if its carriers are pairwise disjoint, some
      possibly empty;
\item a \emph{\(\cup\)-structure} if every carrier is nonempty, the carriers
      possibly overlapping.
\end{itemize}
A \emph{context} is a set \(\Delta\subseteq\mathrm{Var}_{\mathrm{MS}}\).  A
\emph{\(\Delta\)-assignment} into \(\mathcal M\) is a function \(g\) on
\(\Delta\) with \(g(v)\in M_{\mathsf s}\) for every \(v\in\Delta\) of sort
\(\mathsf s\).  For \(v\) of sort \(\mathsf s\) and \(b\in M_{\mathsf s}\),
\(g[v\mapsto b]\) is the \((\Delta\cup\{v\})\)-assignment that sends \(v\) to
\(b\) and agrees with \(g\) elsewhere.  A
\(\mathrm{Var}_{\mathrm{MS}}\)-assignment is a \emph{total} assignment; these
are the assignments of \cref{def:p02-many-sorted-structure}.
\index{many-sorted first-order logic!sorted structure}
\index{assignment!contextual}
\end{hbdefinition}

A sort with no element has no total assignment, because its variable stock is
nonempty.  Contexts are therefore what lets an empty sort carry any
assignment at all.

\begin{hbdefinition}{def:p02-w02-p02-038-satisfaction}{Term values and satisfaction in context}
Let \(g\) be a \(\Delta\)-assignment into a sorted structure \(\mathcal M\).
For a term \(t\) of sort \(\mathsf s\) whose variables lie in \(\Delta\), the
value \(\llbracket t\rrbracket^{\mathcal M}_{g}\in M_{\mathsf s}\) is defined
by recursion on \(t\):
\[
 \llbracket v\rrbracket^{\mathcal M}_{g}\defeq g(v),\qquad
 \llbracket c\rrbracket^{\mathcal M}_{g}\defeq c^{\mathcal M},
\]
\[
 \llbracket f(t_1,\ldots,t_n)\rrbracket^{\mathcal M}_{g}\defeq
 f^{\mathcal M}\bigl(\llbracket t_1\rrbracket^{\mathcal M}_{g},\ldots,
   \llbracket t_n\rrbracket^{\mathcal M}_{g}\bigr).
\]
The value lies in \(M_{\mathsf s}\) because each argument value lies in the
carrier of its declared sort.  For a formula \(\varphi\) with
\(\FV(\varphi)\subseteq\Delta\), the relation
\(\mathcal M\entails\varphi[g]\) is defined by recursion on \(\varphi\):
\[
\begin{aligned}
 \mathcal M\entails P(t_1,\ldots,t_n)[g]
 &\quad\text{iff}\quad
 (\llbracket t_1\rrbracket^{\mathcal M}_{g},\ldots,
   \llbracket t_n\rrbracket^{\mathcal M}_{g})\in P^{\mathcal M},\\
 \mathcal M\entails (t=u)[g]
 &\quad\text{iff}\quad
 \llbracket t\rrbracket^{\mathcal M}_{g}=\llbracket u\rrbracket^{\mathcal M}_{g},\\
 \mathcal M\not\entails\bottom[g],&\\
 \mathcal M\entails\neg\varphi[g]
 &\quad\text{iff}\quad \mathcal M\not\entails\varphi[g],\\
 \mathcal M\entails(\varphi\land\psi)[g]
 &\quad\text{iff}\quad
 \mathcal M\entails\varphi[g]\text{ and }\mathcal M\entails\psi[g],\\
 \mathcal M\entails(\varphi\lor\psi)[g]
 &\quad\text{iff}\quad
 \mathcal M\entails\varphi[g]\text{ or }\mathcal M\entails\psi[g],\\
 \mathcal M\entails(\varphi\to\psi)[g]
 &\quad\text{iff}\quad
 \mathcal M\not\entails\varphi[g]\text{ or }\mathcal M\entails\psi[g],\\
 \mathcal M\entails\forall v\,\varphi[g]
 &\quad\text{iff}\quad
 \mathcal M\entails\varphi[g[v\mapsto b]]\text{ for every }b\in M_{\mathsf s},\\
 \mathcal M\entails\exists v\,\varphi[g]
 &\quad\text{iff}\quad
 \mathcal M\entails\varphi[g[v\mapsto b]]\text{ for some }b\in M_{\mathsf s},
\end{aligned}
\]
where \(v\) has sort \(\mathsf s\).  In the strict language the equality
clause applies only to terms of one sort.  The \emph{cross-sort} language
also admits \(t=u\) for terms of different sorts, with the same clause.  For
a strict structure and a total assignment \(s\), \(\mathcal M\entails\varphi[s]\)
is the relation \(\mathcal M,s\entails\varphi\) of
\cref{def:p02-many-sorted-structure}, whose clauses are the ones displayed.
\index{satisfaction!many-sorted, in context}
\end{hbdefinition}

When a sort \(\mathsf s\) has \(M_{\mathsf s}=\varnothing\), the clauses make
\(\forall v\,\varphi\) true and \(\exists v\,\varphi\) false at every
assignment, for every \(v\) of sort \(\mathsf s\).  This is the source of
every difference between the empty-sort variant and the strict convention in
\cref{sec:p02-w02-p02-038-empty}.

\begin{hblemma}{thm:p02-w02-p02-038-coincidence}{proved=annex:p02-w02-p02-038-coincidence}{Coincidence for contextual assignments}
Let \(\mathcal M\) be a sorted structure, let \(g\) be a \(\Delta\)-assignment
and \(g'\) a \(\Delta'\)-assignment into \(\mathcal M\), and suppose they agree
on a set containing every variable of a term \(t\) or every variable in
\(\FV(\varphi)\).  Then
\(\llbracket t\rrbracket^{\mathcal M}_{g}=\llbracket t\rrbracket^{\mathcal M}_{g'}\),
and \(\mathcal M\entails\varphi[g]\) iff \(\mathcal M\entails\varphi[g']\).
In particular, if \(\FV(\varphi)\subseteq\Delta\subseteq\Delta'\), then
\(\mathcal M\entails\varphi[g']\) iff
\(\mathcal M\entails\varphi[g'{\restriction}\Delta]\).
\end{hblemma}

\subsection{The sort conventions of the reduction, and coverage}
\label{sec:p02-w02-p02-038-coverage}

\begin{hbconvention}{def:p02-w02-p02-038-axiom-families}{The four sort-axiom families, and no coverage axiom}
The theory \(\mathsf{Ax}_{\mathrm{sort}}(\Lang_{\mathrm{MS}})\) of
\cref{def:p02-many-sorted-translation} consists of exactly four families,
named here for reference:
\[
\begin{array}{ll}
 (\mathrm{Ne}) & \exists x\,Q_{\mathsf s}(x)\quad(\mathsf s\in\mathsf{Sort}),\\[2pt]
 (\mathrm{Ct}) & Q_{\mathsf s}(c)\quad(c:\longrightarrow\mathsf s),\\[2pt]
 (\mathrm{Cl}) & \forall x_1\cdots\forall x_n
   \bigl(\bigwedge_{i=1}^{n}Q_{\mathsf s_i}(x_i)\to
   Q_{\mathsf s}(f(x_1,\ldots,x_n))\bigr)\\
   & \qquad(f:\mathsf s_1\times\cdots\times\mathsf s_n\longrightarrow\mathsf s,\ n\geq1),\\[2pt]
 (\mathrm{Dj}) & \forall x\,\neg\bigl(Q_{\mathsf s}(x)\land Q_{\mathsf t}(x)\bigr)
   \quad(\mathsf s\ne\mathsf t).
\end{array}
\]
These express nonemptiness, constant typing, function closure and
disjointness.  \emph{The reduction contains no coverage axiom.}  A one-sorted
model of
\(\mathsf{Ax}_{\mathrm{sort}}(\Lang_{\mathrm{MS}})\cup
\mathsf{VarGuard}_{\mathrm{MS}}(\Lang_{\mathrm{MS}})\) may contain elements
that satisfy no \(Q_{\mathsf s}\), and the one-sorted structure built in the
proof of \cref{thm:p02-many-sorted-translation} always contains one, the junk
point.  Such elements are harmless: a translated formula quantifies only
under a guard \(Q_{\mathsf s}\), and, under an assignment satisfying the
guards, the value of every term lies in some \(Q_{\mathsf s}\)
(\cref{thm:p02-w02-p02-038-restriction}).  Two subsets are used below:
\[
 \mathsf{Ax}^{\mathrm E}_{\mathrm{sort}}\defeq
   (\mathrm{Ct})\cup(\mathrm{Cl})\cup(\mathrm{Dj}),
 \qquad
 \mathsf{Ax}^{\cup}_{\mathrm{sort}}\defeq
   (\mathrm{Ne})\cup(\mathrm{Ct})\cup(\mathrm{Cl}).
\]
For a context \(\Delta\), the \emph{context guard} is
\(\mathsf{Guard}(\Delta)\defeq\{Q_{\mathsf s}(v):v\in\Delta\text{ of sort }
\mathsf s\}\); thus
\(\mathsf{VarGuard}_{\mathrm{MS}}(\Lang_{\mathrm{MS}})=
\mathsf{Guard}(\mathrm{Var}_{\mathrm{MS}})\).  A one-sorted assignment
\(\sigma\) into \(\mathcal A\) satisfies \(\mathsf{Guard}(\Delta)\) exactly
when \(\sigma(v)\in Q_{\mathsf s}^{\mathcal A}\) for every \(v\in\Delta\) of
sort \(\mathsf s\).
\index{sort axiom!coverage}
\end{hbconvention}

The nonemptiness family and the assignment guard both force every
\(Q_{\mathsf s}\) to be inhabited, so in the presence of the guard the first
family adds nothing.

\begin{hblemma}{thm:p02-w02-p02-038-guard-nonempty}{proved=annex:p02-w02-p02-038-guard-nonempty}{The assignment guard implies the nonemptiness family}
If a one-sorted structure \(\mathcal A\) and assignment \(\sigma\) satisfy
\(\mathsf{VarGuard}_{\mathrm{MS}}(\Lang_{\mathrm{MS}})\), then
\(\mathcal A\entails\exists x\,Q_{\mathsf s}(x)\) for every sort
\(\mathsf s\).  Consequently, for every set \(X\) of
\(\Lang_{\mathrm{MS}}^{\flat}\)-formulas and every formula \(\chi\),
\(\mathsf{Ax}_{\mathrm{sort}}\cup\mathsf{VarGuard}_{\mathrm{MS}}\cup
X\entails_{\mathrm{FOL}}\chi\) iff
\(\mathsf{Ax}^{\mathrm E}_{\mathrm{sort}}\cup\mathsf{VarGuard}_{\mathrm{MS}}
\cup X\entails_{\mathrm{FOL}}\chi\).
\end{hblemma}

\begin{hbdefinition}{def:p02-w02-p02-038-coverage}{The coverage variant}
Suppose \(\mathsf{Sort}=\{\mathsf s_1,\ldots,\mathsf s_k\}\) is finite.  The
\emph{coverage axiom} is the one-sorted sentence
\[
 \mathsf{Cov}(\Lang_{\mathrm{MS}})\defeq
 \forall x\,\bigl(Q_{\mathsf s_1}(x)\lor\cdots\lor Q_{\mathsf s_k}(x)\bigr).
\]
The \emph{coverage variant} of the reduction adds
\(\mathsf{Cov}(\Lang_{\mathrm{MS}})\) to the premises of
\cref{thm:p02-many-sorted-translation}.  Finiteness of \(\mathsf{Sort}\) is
a precondition: a first-order formula is a finite string, so the disjunction
exists only for finitely many sorts, and by
\cref{thm:p02-w02-p02-038-coverage-not-elementary} no set of first-order
sentences can take its place when there are infinitely many sorts.
\index{sort axiom!coverage}
\end{hbdefinition}

\begin{hbtheorem}{thm:p02-w02-p02-038-coverage-not-elementary}{proved=annex:p02-w02-p02-038-coverage-not-elementary}{Coverage is not first-order with infinitely many sorts}
Assume the Axiom of Choice in the ambient set theory, and let
\(\mathsf{Sort}\) be infinite.  There is no set \(T\) of
\(\Lang_{\mathrm{MS}}^{\flat}\)-sentences such that, for every one-sorted
model \(\mathcal A\) of \(\mathsf{Ax}_{\mathrm{sort}}(\Lang_{\mathrm{MS}})\),
\(\mathcal A\entails T\) iff every element of \(\mathcal A\) lies in some
\(Q_{\mathsf s}^{\mathcal A}\).
\end{hbtheorem}

\subsection{The satisfaction correspondence in both directions}
\label{sec:p02-w02-p02-038-correspondence}

Two constructions carry the reduction.  The first restricts a one-sorted model
of the closure axioms to the sets cut out by the sort predicates.  The second
transports a sorted structure into a one-sorted domain along injections, one
per sort, and sends every ill-typed function input to a value chosen by a
fixed rule.  The tagged structure of \cref{thm:p02-many-sorted-translation},
its coverage modification, and the plain-union structure used for overlapping
sorts are three instances of the second construction.

\begin{hbdefinition}{def:p02-w02-p02-038-restriction}{Restriction of a one-sorted model}
Let \(\mathcal A\) be an \(\Lang_{\mathrm{MS}}^{\flat}\)-structure satisfying
\((\mathrm{Ct})\) and \((\mathrm{Cl})\).  Its \emph{restriction}
\(\mathcal A^{\downarrow}\) is the sorted structure with
\[
 M_{\mathsf s}\defeq Q_{\mathsf s}^{\mathcal A},\qquad
 c^{\mathcal A^{\downarrow}}\defeq c^{\mathcal A},
\]
\[
 f^{\mathcal A^{\downarrow}}\defeq
   f^{\mathcal A}{\restriction}(M_{\mathsf s_1}\times\cdots\times M_{\mathsf s_n}),
\]
\[
 P^{\mathcal A^{\downarrow}}\defeq
   P^{\mathcal A}\cap(M_{\mathsf s_1}\times\cdots\times M_{\mathsf s_n}).
\]
For a one-sorted assignment \(\sigma\) satisfying \(\mathsf{Guard}(\Delta)\),
\(\sigma{\restriction}\Delta\) is a \(\Delta\)-assignment into
\(\mathcal A^{\downarrow}\).
\end{hbdefinition}

\begin{hbdefinition}{def:p02-w02-p02-038-transport}{Transport along sort injections}
Let \(\mathcal M\) be a sorted structure, \(D\) a nonempty set, and
\(h=(h_{\mathsf s})_{\mathsf s\in\mathsf{Sort}}\) a family of injections
\(h_{\mathsf s}:M_{\mathsf s}\to D\).  For each function symbol \(f\) of
positive arity \(n\) let \(j_f:D^n\to D\) be a given \emph{fill} function.
The one-sorted structure \(\mathcal M^{h}\) has domain \(D\) and
\[
\begin{aligned}
 Q_{\mathsf s}^{\mathcal M^{h}}&\defeq h_{\mathsf s}[M_{\mathsf s}],\qquad
 c^{\mathcal M^{h}}\defeq h_{\mathsf s}(c^{\mathcal M})\quad(c:\longrightarrow\mathsf s),\\
 f^{\mathcal M^{h}}(d_1,\ldots,d_n)&\defeq
 \begin{cases}
 h_{\mathsf s}\bigl(f^{\mathcal M}(b_1,\ldots,b_n)\bigr)
   & \text{if each }d_i=h_{\mathsf s_i}(b_i),\ b_i\in M_{\mathsf s_i},\\
 j_f(d_1,\ldots,d_n) & \text{otherwise},
 \end{cases}\\
 P^{\mathcal M^{h}}&\defeq\bigl\{\bigl(h_{\mathsf s_1}(b_1),\ldots,
   h_{\mathsf s_n}(b_n)\bigr):\\
 &\qquad\qquad(b_1,\ldots,b_n)\in P^{\mathcal M}\bigr\},
\end{aligned}
\]
for \(f:\mathsf s_1\times\cdots\times\mathsf s_n\longrightarrow\mathsf s\)
and \(P:\mathsf s_1\times\cdots\times\mathsf s_n\).  The first case of
\(f^{\mathcal M^{h}}\) is well defined because each \(h_{\mathsf s_i}\) is
injective.  The family \(h\) is \emph{coherent} if
\(h_{\mathsf s}(a)=h_{\mathsf t}(b)\) iff \(a=b\), for all sorts
\(\mathsf s,\mathsf t\), \(a\in M_{\mathsf s}\) and \(b\in M_{\mathsf t}\).
For a \(\Delta\)-assignment \(g\) into \(\mathcal M\), a one-sorted assignment
\(\sigma\) into \(\mathcal M^{h}\) \emph{transports} \(g\) if
\(\sigma(v)=h_{\mathsf s}(g(v))\) for every \(v\in\Delta\) of sort
\(\mathsf s\).  The three instances used below are:
\begin{enumerate}
\item \emph{tagged with junk point}: \(D=\{(\mathsf s,a):\mathsf s\in
      \mathsf{Sort},\ a\in M_{\mathsf s}\}\cup\{\star\}\) with
      \(\star\defeq(\mathsf{Sort},\mathsf{Sort})\),
      \(h_{\mathsf s}(a)=(\mathsf s,a)\), and every \(j_f\) constantly
      \(\star\); the assignment sends variables outside \(\Delta\) to
      \(\star\).  This structure is written \(\mathcal M^{\sharp}\);
\item \emph{tagged with defaults}, for a strict \(\mathcal M\) with a total
      assignment \(s\): the same \(D\) without \(\star\), the same
      \(h\), and \(j_f\) constantly \(h_{\mathsf s}(s(x^{\mathsf s}))\) for
      \(f\) of codomain \(\mathsf s\); one-sorted variables outside the
      image of the variable injection go to one fixed element of \(D\).
      This structure is written \(\mathcal M^{\sharp}_{s}\);
\item \emph{plain union with junk point}: \(D=U\cup\{U\}\) with
      \(U\defeq\bigcup_{\mathsf s}M_{\mathsf s}\), each \(h_{\mathsf s}\) the
      inclusion, every \(j_f\) constantly \(U\), and variables outside
      \(\Delta\) sent to \(U\).  This structure is written
      \(\mathcal M^{\cup}\).
\end{enumerate}
By Foundation, \(\mathsf{Sort}\notin\mathsf{Sort}\) and \(U\notin U\), so
\(\star\) is no tagged pair \((\mathsf s,a)\) with \(\mathsf s\in\mathsf{Sort}\)
and \(U\notin U\); both junk points are fresh without any choice.
\index{many-sorted first-order logic!tagged structure}
\end{hbdefinition}

\begin{hblemma}{thm:p02-w02-p02-038-restriction}{proved=annex:p02-w02-p02-038-restriction}{Restriction lemma}
Let \(\mathcal A\) satisfy \((\mathrm{Ct})\) and \((\mathrm{Cl})\), let
\(\Delta\) be a context, and let \(\sigma\) be a one-sorted assignment into
\(\mathcal A\) satisfying \(\mathsf{Guard}(\Delta)\).  Then:
\begin{enumerate}
\item \(\mathcal A^{\downarrow}\) is a sorted structure; it is an E-structure
      if \(\mathcal A\entails(\mathrm{Dj})\), a \(\cup\)-structure if
      \(\mathcal A\entails(\mathrm{Ne})\), and strict if
      \(\mathcal A\entails\mathsf{Ax}_{\mathrm{sort}}\);
\item for every term \(t\) of sort \(\mathsf s\) with variables in
      \(\Delta\),
      \(\llbracket t\rrbracket^{\mathcal A}_{\sigma}=
      \llbracket t\rrbracket^{\mathcal A^{\downarrow}}_{\sigma\restriction\Delta}
      \in Q_{\mathsf s}^{\mathcal A}\);
\item for every formula \(\varphi\) of the cross-sort language with
      \(\FV(\varphi)\subseteq\Delta\),
      \(\mathcal A,\sigma\entails\operatorname{tr}_{\mathrm{MS}}(\varphi)\)
      iff
      \(\mathcal A^{\downarrow}\entails\varphi[\sigma{\restriction}\Delta]\).
\end{enumerate}
Here \(\operatorname{tr}_{\mathrm{MS}}\) leaves every atom unchanged,
including a cross-sort equality.
\end{hblemma}

\begin{hblemma}{thm:p02-w02-p02-038-transport}{proved=annex:p02-w02-p02-038-transport}{Transport lemma}
Let \(\mathcal M^{h}\) be as in \cref{def:p02-w02-p02-038-transport}, let
\(g\) be a \(\Delta\)-assignment into \(\mathcal M\), and let \(\sigma\)
transport \(g\).  Then:
\begin{enumerate}
\item for every term \(t\) of sort \(\mathsf s\) with variables in
      \(\Delta\),
      \(\llbracket t\rrbracket^{\mathcal M^{h}}_{\sigma}=
      h_{\mathsf s}\bigl(\llbracket t\rrbracket^{\mathcal M}_{g}\bigr)\);
\item for every formula \(\varphi\) of the strict language with
      \(\FV(\varphi)\subseteq\Delta\),
      \(\mathcal M\entails\varphi[g]\) iff
      \(\mathcal M^{h},\sigma\entails\operatorname{tr}_{\mathrm{MS}}(\varphi)\);
      if \(h\) is coherent, the same holds for every formula of the
      cross-sort language;
\item \(\mathcal M^{h}\) satisfies \((\mathrm{Ct})\) and \((\mathrm{Cl})\)
      whatever the fill functions are; it satisfies \((\mathrm{Ne})\) iff
      every \(M_{\mathsf s}\) is nonempty; \((\mathrm{Dj})\) iff the images
      \(h_{\mathsf s}[M_{\mathsf s}]\) are pairwise disjoint; and, for
      finite \(\mathsf{Sort}\), \(\mathsf{Cov}\) iff these images cover
      \(D\);
\item \(\sigma\) satisfies \(\mathsf{Guard}(\Delta)\).
\end{enumerate}
\end{hblemma}

The three instances have the following properties, each checked in
\cref{annex:p02-w02-p02-038-transport}.  The tagged images are always
pairwise disjoint, and the tagged family is coherent iff \(\mathcal M\) has
pairwise disjoint carriers.  The plain-union family is always coherent, and
its images are disjoint iff \(\mathcal M\) has pairwise disjoint carriers.
The junk point lies in no \(Q_{\mathsf s}\), so, when \(\mathsf{Sort}\) is
finite, neither junk structure satisfies \(\mathsf{Cov}\); the tagged
structure with defaults always does.

\begin{hbtheorem}{thm:p02-w02-p02-038-strict-correspondence}{proved=annex:p02-w02-p02-038-strict-correspondence}{Two-way satisfaction correspondence for the strict reduction}
Let \(\Lang_{\mathrm{MS}}\) be a strict many-sorted signature.
\begin{enumerate}
\item \emph{Restriction.}  If \(\mathcal A\entails
      \mathsf{Ax}_{\mathrm{sort}}(\Lang_{\mathrm{MS}})\) and \(\sigma\)
      satisfies \(\mathsf{VarGuard}_{\mathrm{MS}}(\Lang_{\mathrm{MS}})\),
      then \(\mathcal A^{\downarrow}\) is strict,
      \(\sigma{\restriction}\mathrm{Var}_{\mathrm{MS}}\) is a total
      assignment into it, and for every well-sorted \(\varphi\),
      \(\mathcal A,\sigma\entails\operatorname{tr}_{\mathrm{MS}}(\varphi)\)
      iff
      \(\mathcal A^{\downarrow},\sigma{\restriction}\mathrm{Var}_{\mathrm{MS}}
      \entails\varphi\).
\item \emph{Tagging.}  If \(\mathcal M\) is strict and \(s\) is a total
      assignment, let \(s^{\sharp}\) send each sorted variable \(v\) of sort
      \(\mathsf s\) to \((\mathsf s,s(v))\) and every other one-sorted
      variable to \(\star\).  Then \(\mathcal M^{\sharp}\entails
      \mathsf{Ax}_{\mathrm{sort}}(\Lang_{\mathrm{MS}})\), \(s^{\sharp}\)
      satisfies \(\mathsf{VarGuard}_{\mathrm{MS}}(\Lang_{\mathrm{MS}})\), and
      for every well-sorted \(\varphi\),
      \(\mathcal M,s\entails\varphi\) iff
      \(\mathcal M^{\sharp},s^{\sharp}\entails
      \operatorname{tr}_{\mathrm{MS}}(\varphi)\).
\item \emph{Consequence.}  For every set \(\Gamma\cup\{\varphi\}\) of
      well-sorted formulas, the equivalence of
      \cref{thm:p02-many-sorted-translation} holds.
\end{enumerate}
\end{hbtheorem}

Clause~3 is the reduction theorem of \cref{sec:p02-many-sorted}; the annex
proof of this theorem writes out the induction cases and the axiom checks
that the proof attached to \cref{thm:p02-many-sorted-translation} summarizes.

\begin{hbtheorem}{thm:p02-w02-p02-038-coverage-reduction}{proved=annex:p02-w02-p02-038-coverage-reduction}{The coverage variant of the reduction}
Let \(\mathsf{Sort}\) be finite, and let \(\Gamma\cup\{\varphi\}\) be a set
of well-sorted formulas.  Then
\[
 \Gamma\entails_{\mathrm{MS}}\varphi
 \quad\Longleftrightarrow\quad
 \mathsf{Ax}_{\mathrm{sort}}\cup\{\mathsf{Cov}\}
   \cup\mathsf{VarGuard}_{\mathrm{MS}}
   \cup\operatorname{tr}_{\mathrm{MS}}[\Gamma]
 \entails_{\mathrm{FOL}}\operatorname{tr}_{\mathrm{MS}}(\varphi).
\]
Hence adding \(\mathsf{Cov}\) does not change which translated formulas
follow from translated premises.  In the converse construction the junk point
is removed and every ill-typed input of a function symbol of codomain
\(\mathsf s\) is sent to the tagged value of the designated variable
\(x^{\mathsf s}\) under the given assignment.  The structure
\(\mathcal M^{\sharp}\) of \cref{thm:p02-w02-p02-038-strict-correspondence}
cannot be used, since \(\star\) violates \(\mathsf{Cov}\).
\end{hbtheorem}

\begin{hbexample}{def:p02-w02-p02-038-example}{Points and lines}
Let \(\mathsf{Sort}=\{\mathsf p,\mathsf l\}\), with a relation symbol
\(I:\mathsf p\times\mathsf l\) (incidence), a constant \(o:\longrightarrow
\mathsf p\), and a function symbol \(m:\mathsf l\longrightarrow\mathsf p\)
(a marked point on a line).  The well-sorted sentence
\(\varphi\defeq\forall y^{\mathsf l}\,I(m(y^{\mathsf l}),y^{\mathsf l})\)
translates to
\(\forall y^{\mathsf l}\bigl(Q_{\mathsf l}(y^{\mathsf l})\to
I(m(y^{\mathsf l}),y^{\mathsf l})\bigr)\), and
\(\mathsf{Ax}_{\mathrm{sort}}\) consists of \(\exists x\,Q_{\mathsf p}(x)\),
\(\exists x\,Q_{\mathsf l}(x)\), \(Q_{\mathsf p}(o)\),
\(\forall x\,(Q_{\mathsf l}(x)\to Q_{\mathsf p}(m(x)))\) and the two
disjointness axioms for the ordered pairs of distinct sorts.

Take the strict structure \(M_{\mathsf p}=\{a\}\),
\(M_{\mathsf l}=\{\ell\}\), \(o^{\mathcal M}=a\), \(m^{\mathcal M}(\ell)=a\),
\(I^{\mathcal M}=\{(a,\ell)\}\).  Then \(\mathcal M\entails\varphi\).  The
tagged structure has domain
\(\{(\mathsf p,a),(\mathsf l,\ell),\star\}\), with
\[
 Q_{\mathsf p}=\{(\mathsf p,a)\},\quad Q_{\mathsf l}=\{(\mathsf l,\ell)\},\quad
 o=(\mathsf p,a),\quad
 I=\{((\mathsf p,a),(\mathsf l,\ell))\},
\]
and
\[
 m^{\mathcal M^{\sharp}}((\mathsf l,\ell))=(\mathsf p,a),\qquad
 m^{\mathcal M^{\sharp}}((\mathsf p,a))=\star,\qquad
 m^{\mathcal M^{\sharp}}(\star)=\star.
\]
The closure axiom for \(m\) holds because only \((\mathsf l,\ell)\)
satisfies its antecedent; the two junk values occur where the antecedent
fails.  The translated sentence holds, since its guard admits only
\((\mathsf l,\ell)\).  The coverage axiom
\(\forall x\,(Q_{\mathsf p}(x)\lor Q_{\mathsf l}(x))\) fails at \(\star\).
In the structure with defaults the domain is
\(\{(\mathsf p,a),(\mathsf l,\ell)\}\) and
\(m((\mathsf p,a))=(\mathsf p,s(x^{\mathsf p}))=(\mathsf p,a)\); coverage
holds and the translated sentence is still true.

If instead \(M_{\mathsf l}\) were empty, which the strict convention forbids
and the empty-sort variant below allows, \(\varphi\) would be vacuously true
and \(\exists y^{\mathsf l}\,(y^{\mathsf l}=y^{\mathsf l})\) false.
\end{hbexample}

\subsection{A variant admitting empty sorts}
\label{sec:p02-w02-p02-038-empty}

\begin{hbdefinition}{def:p02-w02-p02-038-empty-semantics}{Empty-sort consequence}
Let \(\Gamma\cup\{\varphi\}\) be a set of formulas of the strict language and
\(\Delta\) a context with \(\FV(\gamma)\subseteq\Delta\) for every
\(\gamma\in\Gamma\cup\{\varphi\}\).  Write
\(\Gamma\entails^{\Delta}_{\mathrm E}\varphi\) iff, for every E-structure
\(\mathcal M\) and every \(\Delta\)-assignment \(g\) into \(\mathcal M\), if
\(\mathcal M\entails\gamma[g]\) for every \(\gamma\in\Gamma\) then
\(\mathcal M\entails\varphi[g]\).  Write
\(\Gamma\entails^{\mathrm{tot}}_{\mathrm E}\varphi\) for the same definition
with total assignments in place of \(\Delta\)-assignments.  A sort
\(\mathsf s\) is \emph{inhabited by the signature} if some term of sort
\(\mathsf s\) contains no variable.
\index{many-sorted first-order logic!empty sorts}
\end{hbdefinition}

The handling of assignments is the substance of this variant.  If
assignments stay total, an empty sort admits none, and such structures drop
out of the consequence relation altogether
(\cref{thm:p02-w02-p02-038-total-collapse}).  Contexts let a structure with
an empty sort count, at contexts that contain no variable of that sort.
Consequence then depends on the context, which is the phenomenon Pitts
illustrates for algebraic theories with a sort that may be empty
\cite[Example~2.3.2, p.~11]{Pitts1995Draft}.  The contextual powerset
interpretation of \cref{thm:p10-doctrine-tarski} keeps the context in the
semantic judgment for the same reason.

\begin{hblemma}{thm:p02-w02-p02-038-total-collapse}{proved=annex:p02-w02-p02-038-total-collapse}{Total assignments hide empty sorts}
Every E-structure that has a total assignment is strict.  Consequently
\(\Gamma\entails^{\mathrm{tot}}_{\mathrm E}\varphi\) iff
\(\Gamma\entails_{\mathrm{MS}}\varphi\).
\end{hblemma}

\begin{hblemma}{thm:p02-w02-p02-038-inhabited}{proved=annex:p02-w02-p02-038-inhabited}{Sorts that can be empty}
For each sort \(\mathsf s\), the following are equivalent: \(\mathsf s\) is
not inhabited by the signature; some E-structure has
\(M_{\mathsf s}=\varnothing\).  Moreover there is one E-structure
\(\mathcal M_{0}\) in which exactly the sorts not inhabited by the signature
are empty.
\end{hblemma}

\begin{hbtheorem}{thm:p02-w02-p02-038-empty-reduction}{proved=annex:p02-w02-p02-038-empty-reduction}{Reduction for the empty-sort variant}
Let \(\Gamma\cup\{\varphi\}\) be formulas of the strict language and
\(\Delta\) a context containing all their free variables.  Then
\[
 \Gamma\entails^{\Delta}_{\mathrm E}\varphi
 \quad\Longleftrightarrow\quad
 \mathsf{Ax}^{\mathrm E}_{\mathrm{sort}}\cup\mathsf{Guard}(\Delta)
   \cup\operatorname{tr}_{\mathrm{MS}}[\Gamma]
 \entails_{\mathrm{FOL}}\operatorname{tr}_{\mathrm{MS}}(\varphi).
\]
\end{hbtheorem}

Compared with \cref{thm:p02-many-sorted-translation}, three things change.
The nonemptiness family is dropped: it would make
\(\operatorname{tr}_{\mathrm{MS}}(\exists x^{\mathsf s}\,
(x^{\mathsf s}=x^{\mathsf s}))\) a consequence of the sort axioms alone.  The
total guard \(\mathsf{VarGuard}_{\mathrm{MS}}\) is replaced by the context
guard \(\mathsf{Guard}(\Delta)\): by
\cref{thm:p02-w02-p02-038-guard-nonempty}, the total guard would reinstate
nonemptiness of every sort.  And the junk point is no longer only a home for
ill-typed function values: when every carrier is empty, it is the only
element of the one-sorted domain, which ordinary first-order semantics
requires to be nonempty.

\begin{hbtheorem}{thm:p02-w02-p02-038-empty-comparison}{proved=annex:p02-w02-p02-038-empty-comparison}{Empty-sort consequence compared with the strict convention}
Let \(\Gamma\cup\{\varphi\}\) be formulas of the strict language and
\(\Delta\) a context containing all their free variables.
\begin{enumerate}
\item If \(\Gamma\entails^{\Delta}_{\mathrm E}\varphi\), then
      \(\Gamma\entails_{\mathrm{MS}}\varphi\).
\item If \(\mathsf s\) is not inhabited by the signature, \(v\) is a variable
      of sort \(\mathsf s\), and \(\FV(\psi)\subseteq\{v\}\), then
      \(\forall v\,\psi\not\entails^{\varnothing}_{\mathrm E}\exists v\,\psi\),
      although \(\forall v\,\psi\entails_{\mathrm{MS}}\exists v\,\psi\).
      In particular \(\exists v\,(v=v)\) is valid in the strict sense but
      not in the empty-sort sense.
\item Enlarging the context is sound: if \(\Delta\subseteq\Delta'\) then
      \(\Gamma\entails^{\Delta}_{\mathrm E}\varphi\) implies
      \(\Gamma\entails^{\Delta'}_{\mathrm E}\varphi\).  Shrinking it is not:
      for \(\mathsf s\) and \(v\) as in clause~2 and a further variable
      \(w\) of sort \(\mathsf s\),
      \(\varnothing\entails^{\{w\}}_{\mathrm E}\exists v\,(v=v)\) but
      \(\varnothing\not\entails^{\varnothing}_{\mathrm E}\exists v\,(v=v)\).
\item If every sort is inhabited by the signature, then every E-structure
      is strict and \(\Gamma\entails^{\Delta}_{\mathrm E}\varphi\) iff
      \(\Gamma\entails_{\mathrm{MS}}\varphi\).  This clause uses no choice:
      it fixes one injection of the at most countable symbol set into the
      natural numbers.
\item Over strict structures only, contextual and total consequence agree:
      \(\Gamma\entails_{\mathrm{MS}}\varphi\) iff the condition defining
      \(\Gamma\entails^{\Delta}_{\mathrm E}\varphi\) holds for every strict
      \(\mathcal M\) and every \(\Delta\)-assignment.  The right-to-left
      half uses no choice.  The left-to-right half extends a
      \(\Delta\)-assignment to a total one with a choice function on the
      family of carriers; it needs the Axiom of Choice only when
      \(\mathsf{Sort}\) is infinite.
\end{enumerate}
\end{hbtheorem}

\begin{hbwarning}{def:p02-w02-p02-038-empty-warning}{Rules that need a nonempty sort}
Clause~2 of \cref{thm:p02-w02-p02-038-empty-comparison} shows that the
passage from \(\forall v\,\psi\) to \(\exists v\,\psi\), and every rule that
yields it, such as instantiating a universal at a variable of sort
\(\mathsf s\) that does not occur in the context and then generalizing
existentially, is unsound under the empty-sort convention.  A calculus
for that convention must record contexts and remove a variable from a
context only when its sort is inhabited by the signature
\cite[Example~2.3.2, p.~11]{Pitts1995Draft}.  An empty sort is also
different from a nondenoting term: free logics
(\cref{sec:p02-free-logic}) keep a nonempty inner domain and let a term
fail to denote, whereas here every term still denotes and a sort may have no
elements.
\end{hbwarning}

\subsection{A variant admitting overlapping sorts}
\label{sec:p02-w02-p02-038-overlap}

\Cref{sec:p02-many-sorted} adopts disjoint carriers and notes that
overlapping carriers require a different reconstruction.  The reconstruction
is as follows.

\begin{hbdefinition}{def:p02-w02-p02-038-overlap-semantics}{Overlapping-sort semantics}
The \emph{overlapping-sort} presentation keeps the signature of
\cref{def:p02-many-sorted-signature}, with one declaration for each symbol,
and uses the cross-sort language of \cref{def:p02-w02-p02-038-satisfaction}:
an equation \(t=u\) is a formula for terms of any two sorts, and is true iff
the two values are the same element.  Structures are the
\(\cup\)-structures of \cref{def:p02-w02-p02-038-sorted-structure}: every
carrier is nonempty, and carriers may share elements.  Write
\(\Gamma\entails_{\cup}\varphi\) for consequence over all
\(\cup\)-structures and total assignments.  The guard translation
\(\operatorname{tr}_{\mathrm{MS}}\) is extended by leaving cross-sort
equalities unchanged.
\index{many-sorted first-order logic!overlapping sorts}
\end{hbdefinition}

\begin{hbtheorem}{thm:p02-w02-p02-038-overlap-reduction}{proved=annex:p02-w02-p02-038-overlap-reduction}{One-sorted reduction for overlapping sorts}
For every set \(\Gamma\cup\{\varphi\}\) of formulas of the cross-sort
language,
\[
 \Gamma\entails_{\cup}\varphi
 \quad\Longleftrightarrow\quad
 \mathsf{Ax}^{\cup}_{\mathrm{sort}}\cup\mathsf{VarGuard}_{\mathrm{MS}}
   \cup\operatorname{tr}_{\mathrm{MS}}[\Gamma]
 \entails_{\mathrm{FOL}}\operatorname{tr}_{\mathrm{MS}}(\varphi).
\]
The converse construction is the plain union \(\mathcal M^{\cup}\), not the
tagged union.
\end{hbtheorem}

\begin{hbtheorem}{thm:p02-w02-p02-038-overlap-comparison}{proved=annex:p02-w02-p02-038-overlap-comparison}{Overlapping sorts compared with the strict convention}
Let \(\mathsf s\ne\mathsf t\) be sorts.
\begin{enumerate}
\item For formulas of the strict language, \(\Gamma\entails_{\cup}\varphi\)
      iff \(\Gamma\entails_{\mathrm{MS}}\varphi\).
\item The cross-sort sentence
      \(\theta\defeq\exists x^{\mathsf s}\,\exists x^{\mathsf t}\,
      (x^{\mathsf s}=x^{\mathsf t})\) has a \(\cup\)-model, while
      \(\mathsf{Ax}_{\mathrm{sort}}\entails_{\mathrm{FOL}}
      \neg\operatorname{tr}_{\mathrm{MS}}(\theta)\).  So the disjointness
      family must be dropped from the reduction for overlapping sorts.
\item There is a \(\cup\)-structure \(\mathcal M\) with
      \(\mathcal M\entails\theta\) and
      \(\mathcal M^{\sharp}\not\entails\operatorname{tr}_{\mathrm{MS}}(\theta)\):
      the tagged construction separates an element shared by two sorts and
      so fails for cross-sort equality.
\end{enumerate}
\end{hbtheorem}

Clause~1 says that disjointness is not visible in the strict language: with
only same-sort equality, no formula distinguishes a structure from its tagged
copy.  The difference between the two conventions appears exactly when
cross-sort equality is part of the language.  This presentation does not treat a sort
order with declared inclusions or a symbol with several declarations; each
would add conditions on how the interpretations at different declarations
fit together, and those conditions are not fixed here.

\begin{longtable}{@{}p{.22\linewidth}p{.22\linewidth}p{.22\linewidth}p{.24\linewidth}@{}}
\caption{The three conventions and their one-sorted reductions.}
\label{tab:p02-w02-p02-038-conventions}\\
\toprule
Convention & Carriers and language & One-sorted premises & Converse construction \\ \midrule
\endfirsthead
\toprule
Convention & Carriers and language & One-sorted premises & Converse construction \\ \midrule
\endhead
Strict (\cref{thm:p02-many-sorted-translation}) &
nonempty, pairwise disjoint; same-sort equality; total assignments &
\((\mathrm{Ne})\), \((\mathrm{Ct})\), \((\mathrm{Cl})\), \((\mathrm{Dj})\),
\(\mathsf{VarGuard}_{\mathrm{MS}}\) &
tagged union with junk point \\
\addlinespace
Strict with coverage, finitely many sorts
(\cref{thm:p02-w02-p02-038-coverage-reduction}) &
as strict &
as strict, plus \(\mathsf{Cov}\) &
tagged union, defaults from the assignment \\
\addlinespace
Empty sorts (\cref{thm:p02-w02-p02-038-empty-reduction}) &
possibly empty, pairwise disjoint; same-sort equality; context \(\Delta\) &
\((\mathrm{Ct})\), \((\mathrm{Cl})\), \((\mathrm{Dj})\),
\(\mathsf{Guard}(\Delta)\) &
tagged union with junk point, which may be the only element \\
\addlinespace
Overlapping sorts (\cref{thm:p02-w02-p02-038-overlap-reduction}) &
nonempty, possibly overlapping; cross-sort equality; total assignments &
\((\mathrm{Ne})\), \((\mathrm{Ct})\), \((\mathrm{Cl})\),
\(\mathsf{VarGuard}_{\mathrm{MS}}\) &
plain union with junk point \\
\bottomrule
\end{longtable}

\emph{Relationships.}  The strict correspondence supplies the step-by-step
proof behind \cref{thm:p02-many-sorted-translation}, and through it the
semantic transfers stated after that theorem.  The same assignment guard
\(\mathsf{VarGuard}_{\mathrm{MS}}\) is used by the matched second-order
calculus of \cref{sec:p02-sol-matched-calculus}.  Empty sorts with
contexts are the semantic setting of the contextual first-order logic of
Part~10 (\cref{thm:p10-doctrine-tarski}).  No proof calculus for any of the
variants is fixed here; each reduction concerns semantic consequence only.
